\section{Problem Setup}

\subsection{Model-Heterogeneous CLE}

We consider clients \(k\in\{1,\ldots,K\}\), each with a private dataset and classifier \(f_k\); architectures need not match across clients. The primary controlled setup uses four heterogeneous clients with ResNet10, ResNet12, ShuffleNet, and MobileNetV2 on a CIFAR-10 private task. Client label distributions are non-IID under a Dirichlet partition with concentration \(\alpha=0.5\).

Let \(Y\) denote the task class, \(O\) a concrete corruption operator, and \(F=\pi(O)\) its coarse corruption family. CLE-HFL introduces client-specific class dependence,

\[
P_k(O=o\mid Y=c)\neq P_k(O=o),
\]

with a pre-registered client-specific binding

\[
b_k:\mathcal C\rightarrow\mathcal O.
\]

For \(M=|\mathcal O_{\mathrm{seen}}|\) seen operators, the training operator is sampled according to

\[
P_k(O=o\mid Y=c)
=\frac{1-\gamma_{\mathrm{CLE}}}{M}
+\gamma_{\mathrm{CLE}}\,\mathbf 1\{o=b_k(c)\}.
\]

The control setting uses \(\gamma_{\mathrm{CLE}}=0\), while strong CLE uses \(\gamma_{\mathrm{CLE}}=0.9\), making \(b_k(c)\) a strong but non-deterministic dominant operator for class \(c\). The formal paired evaluation uses 1,000 semantic sources, 15 operators, four clients, and corruption severity 3.

A crucial distinction is granularity. \textbf{CLE-v2 data generation and DSA evaluation are operator-level. The final PEW+BER method is not.} PEW predicts only a coarse corruption-family proxy and BER groups private fit samples by class and that coarse proxy. Operator metadata remain unavailable to the deployable training method. Operator-level labels are used only for CLE construction, sealed DSA evaluation, and non-deployable Oracle boundary audits.

\subsection{Information Separation}

Private data are separated by role. The \textbf{private fit} subset produces local gradients and BER group counts. A \textbf{private audit} subset may support strict communication routing but does not produce local-training gradients. The \textbf{final task test} is used only for frozen utility reporting. The \textbf{paired operator grid} is used only after training for DSA. Public data support PEW training or pre-existing communication without private-task ground-truth leakage.

The train-time method does not read private true operator IDs, true corruption families, evaluation severity, seen/unseen flags, or the paired evaluation binding for optimization, routing, hyperparameter tuning, or model selection. DSA reads the frozen binding only after model sealing. This isolation is required because the diagnostic is deliberately target-aligned.

\section{Paired Diagnosis and Identification}

\subsection{Directional Shortcut Alignment}

For source \(z=(x,y)\), let \(T_o(x)\) be a semantic-preserving corruption transformation. For client \(k\) and operator \(o\), define the classes bound to that operator during training:

\[
B_{k,o}=\{c:b_k(c)=o\}.
\]

We evaluate only source--operator pairs for which \(y\notin B_{k,o}\), avoiding a trivial contribution from the source's true class. Define bound probability mass

\[
m_{k,o}(z,o')=\sum_{c\in B_{k,o}}p_k(c\mid T_{o'}(x)).
\]

The source-level directional contrast is

\[
d_{k,o}(z)=m_{k,o}(z,o)
-\frac{1}{|\mathcal O|-1}\sum_{o'\neq o}m_{k,o}(z,o'),
\]

and

\[
\operatorname{DSA}_{k,o}=\mathbb E_z[d_{k,o}(z)].
\]

Pooled DSA equally averages valid operator-level values and then client-level values:

\[
\operatorname{DSA}_{\mathrm{pool}}
=\frac1K\sum_{k=1}^{K}
\frac1{|\mathcal O_k^{\mathrm{valid}}|}
\sum_{o\in\mathcal O_k^{\mathrm{valid}}}\operatorname{DSA}_{k,o}.
\]

DSA lies on a probability scale: a value of \(0.12\) means roughly a 12-percentage-point directional difference in predictive mass, not a \(0.12\)-point accuracy difference.

The interpretation is narrow: \emph{for the same semantic source, does applying operator \(o\) move probability toward the classes that were associated with \(o\) during this client's training, relative to applying other operators?} The source-paired intervention is a controlled diagnostic under stated assumptions, not an unconditional recovery of the model's full causal mechanism.

\subsection{Paired Cancellation and Binding Specificity}

The following is an identification decomposition of the paired estimand, not a generative assumption about the neural network. Write

\[
m_{k,o}(z,o')=s_{k,o}(z)+r_{k,o}(z,o'),
\]

where \(s_{k,o}(z)\) is an operator-invariant source term and \(r_{k,o}\) is an operator-dependent response. Then

\[
d_{k,o}(z)=r_{k,o}(z,o)
-\frac{1}{|\mathcal O|-1}\sum_{o'\neq o}r_{k,o}(z,o'),
\]

so the operator-invariant source term cancels exactly. Under this decomposition, DSA identifies a \textbf{binding-direction operator-response contrast}. It does not prove that the model uses only corruption, that every operator preserves all task semantics, or that the same response is harmful in every deployment distribution.

To test binding specificity, we preserve the number of classes assigned to each operator but randomly permute the binding and recompute DSA on the same prediction cache. Under the primary strong-CLE HFL checkpoint, true-binding DSA is \(0.119644\), whereas the shuffled-binding null has p95 \(=0.029559\) with permutation \(p=0.000999\). This provides evidence against the explanation that the effect is driven only by generic operator-to-class response patterns unrelated to the training association.

\subsection{Why Consistency and Proxy Scores Are Insufficient}

Within-operator consistency and cross-operator directional invariance are different objects. If multiple predictive views around the same corrupted observation are identical, their JSD is zero; predictions can nevertheless differ systematically across corruption operators. In a fixed-cache construction, within-operator JSD is zero while DSA remains \(0.119644\). Thus,

\[
\begin{aligned}
&\text{within-operator consistency}\\
&\qquad\not\Rightarrow\qquad
\text{cross-operator directional invariance}.
\end{aligned}
\]

This does not make JSD ineffective for corruption robustness; it shows that consistency and DSA answer different questions.

The same separation motivates a proxy non-identifiability result. Let \(Z\) be a proxy environment predicted by PEW and \(E\) the latent true corruption environment. Define

\[
\mathcal D(A,B)=TV(P_{A,B},P_AP_B).
\]

With no constraint on the error channel \(P(Z\mid E,Y)\), a fixed observable \(P(Y,Z)\) can correspond to one latent world with \(\mathcal D(Y,E)=0\) and another with

\[
\mathcal D(Y,E)=1-\sum_eP(E=e)^2>0.
\]

Therefore \(P(Y,Z)\) alone does not identify true environment dependence. A conditional bridge remains valid under sample-wise coupling:

\[
\mathcal D_Q(Y,E)\le
\mathcal D_Q(Y,Z)+2\epsilon_Q,
\qquad
\epsilon_Q=P_Q(Z\neq E).
\]

In the current BER effective distribution, the PEW family error is \(0.507\)--\(0.575\), making this bound trivial after truncation at 1.0. We therefore do not claim an unconditional theoretical guarantee that BER makes the true environment independent of class. Under our evidence protocol, mitigation claims require an additional target-aligned behavioral endpoint; we use paired DSA. Full constructive proofs and algebraic validations are deferred to Appendix A.

\subsection{Statistical Unit}

The independent evaluation unit is the \textbf{semantic source}, not each corrupted image. Multiple operator views of one source share content and must not be treated as independent observations. Main DSA intervals use source-clustered bootstrap after computing source-level contrasts. These intervals condition on the trained checkpoint and quantify evaluation-source uncertainty only; they do \textbf{not} quantify training-seed, partition, binding-map, model-selection, or cross-dataset uncertainty.

\section{Local-versus-Communication Formation Analysis}

We compare four matched arms under CLE-v2: HFL with \(\gamma_{\mathrm{CLE}}=0\), HFL with \(\gamma_{\mathrm{CLE}}=0.9\), Local with \(\gamma_{\mathrm{CLE}}=0\), and Local with \(\gamma_{\mathrm{CLE}}=0.9\). All arms share the same AugMix/JSD/DCL local baseline, use training seed 0, run for 12 rounds, and use at most 16 local batches per client per round. This 12-round budget was frozen for the mechanism-focused factorial and is distinct from the 40-round held-out method comparison; it does not by itself establish long-horizon stability.

\begin{table*}[t]
\centering
\small
\caption{CLE formation and matched local-first attribution. Source-bootstrap intervals condition on the trained checkpoints.}
\label{tab:auto1}
\resizebox{\textwidth}{!}{%
\begin{tabular}{lrrrr}
\toprule
Arm & Scope & \(\gamma_{\mathrm{CLE}}\) & Operator-grid Acc. (\%) & Pooled DSA \\
\midrule
HFL control & HFL & 0.0 & 24.9300 & -0.000245 \\
HFL strong CLE & HFL & 0.9 & 21.4367 & 0.119644 \\
Local control & Local & 0.0 & 25.5683 & -0.001914 \\
Local strong CLE & Local & 0.9 & 22.4233 & 0.106046 \\
\bottomrule
\end{tabular}%
}
\end{table*}

The matched contrasts are

\[
\Delta_{\mathrm{HFL}}=0.119889,\qquad
\Delta_{\mathrm{Local}}=0.107960,
\]

and

\[
\Delta_{\mathrm{comm}}
=\Delta_{\mathrm{HFL}}-\Delta_{\mathrm{Local}}
=0.011929.
\]

Their source-bootstrap 95\% intervals are \([0.118050,0.121562]\), \([0.106145,0.109678]\), and \([0.011115,0.012712]\), respectively. At this evaluated seed and fixed scenario, the descriptive ratio

\[
\Delta_{\mathrm{Local}}/\Delta_{\mathrm{HFL}}=90.05\%
\]

shows that the Local contrast is close to the HFL contrast. We call this \textbf{predominantly local-first formation}. The ratio is not a per-example causal mediation share, and the result does not imply that communication is irrelevant. It motivates targeting the local objective because the matched Local effect already accounts for most of the pooled HFL effect in this setting.

\textbf{[EVIDENCE NEEDED: multi-training-seed stability for the matched HFL-versus-Local four-arm factorial. Candidate S1/S2 repetitions would jointly cover HFL \(\gamma=0\) vs. \(.9\), Local \(\gamma=0\) vs. \(.9\), and local-first stability; these candidate runs are not authorized or reported here.]}

\section{Taxonomy-Assisted Local Mitigation}

\subsection{Public Environment Witness}

The \textbf{Public Environment Witness (PEW)} is trained only on public carrier images with programmatically generated corruption supervision. It is frozen as a \textbf{coarse family-level witness} over

\[
\{\text{clean},\text{noise},\text{blur},\text{weather},\text{digital},\text{unknown}\}.
\]

For each private fit sample, PEW produces a hard pseudo-environment \(\hat e_i\). PEW does not read private true operators or private true corruption families during training. It is therefore private-environment-metadata-free at deployment but explicitly \textbf{taxonomy-assisted}: its semantics come from a predefined corruption-family taxonomy. Architecture and optimizer details are deferred to Appendix B because PEW is treated as side-information infrastructure rather than the paper's architectural contribution.

The granularity mismatch is deliberate. CLE-v2 construction and DSA evaluation remain operator-level, while PEW and BER use only coarse family-level pseudo-environments. The Oracle operator analysis in Appendix C is a non-deployable boundary audit and does not define or motivate an operator-level PEW.

\subsection{Balanced Environment Risk}

For client \(k\), task class \(c\), and pseudo-environment \(e\), define

\[
n_{k,c,e}=\#\{i:y_i=c,\hat e_i=e\},
\qquad
R_{k,c,e}=\frac1{n_{k,c,e}}\sum_{i:y_i=c,\hat e_i=e}\ell_i.
\]

BER assigns class-conditional environment weight

\[
a^{(\gamma_{\mathrm{BER}})}_{k,c,e}
=
\frac{
\mathbf 1[n_{k,c,e}\ge m]
\min(n_{k,c,e},K_{\mathrm{cap}})^{\gamma_{\mathrm{BER}}}
}{
\sum_{e'}
\mathbf 1[n_{k,c,e'}\ge m]
\min(n_{k,c,e'},K_{\mathrm{cap}})^{\gamma_{\mathrm{BER}}}
},
\]

with the frozen Formal configuration

\[
\gamma_{\mathrm{BER}}=0.5,\qquad K_{\mathrm{cap}}=32,\qquad m=2.
\]

The client task risk is

\[
R_{\mathrm{BER},k}
=
\frac1{|\mathcal C_k^{\mathrm{valid}}|}
\sum_{c\in\mathcal C_k^{\mathrm{valid}}}
\sum_e a^{(\gamma_{\mathrm{BER}})}_{k,c,e}R_{k,c,e}.
\]

BER does not directly optimize the worst environment. It reduces the effective risk-mass advantage of large pseudo-environment supports \emph{within each class}. BER does not equalize environment groups; it sublinearly compresses their support-induced risk-mass advantage. Equivalently, it performs empirical risk minimization under effective sample mass

\[
q_i=
\frac{a^{(\gamma_{\mathrm{BER}})}_{k,y_i,\hat e_i}}
{|\mathcal C_k^{\mathrm{valid}}|n_{k,y_i,\hat e_i}}.
\]

For two valid environments within a class,

\[
\frac{Q_\gamma(e_1\mid c)}{Q_\gamma(e_2\mid c)}
=
\left(
\frac{\min(n_{c,e_1},K_{\mathrm{cap}})}
{\min(n_{c,e_2},K_{\mathrm{cap}})}
\right)^{\gamma_{\mathrm{BER}}}.
\]

Hence \(0\le\gamma_{\mathrm{BER}}<1\) compresses the class-conditional log-support advantage. At the current configuration, the maximum effective mass ratio between valid pseudo-environment groups is \(4\). A fixed-data audit verifies that the implemented objective matches this effective-distribution formula; detailed audit statistics are in Appendix B. These statements characterize the \emph{training distribution induced by BER} and do not imply that DSA must decrease or that task accuracy must improve.
